\documentclass[../main.tex]{subfiles}

\begin{document}

\section{Introduction}

The northern permafrost region contains more organic carbon in its upper three meters of soil than currently resides in Earth's atmosphere \citep{hugelius2014estimated}. Typical, gradual permafrost thaw releases stored carbon over decades or longer, while abrupt thaw processes can mobilize deeper carbon stocks in days to years \citep{turetsky2020carbon, nitzbon2020fast}. Under high emissions scenarios, abrupt thaw is projected to impact only 2.5 million km$^2$ out of the 18 million km$^2$ permafrost region, but threatens to release nearly as much carbon as gradual thaw processes operating over a much larger area \citep{turetsky2020carbon}. Abrupt thaw leads to a variety of landforms, including thermokarst lakes, collapsing peatlands, and hillslope failures, all of which differ in the quantity and form of carbon released \citep{olefeldt2016circumpolar, schadel2016potential}. The scale and timing of carbon release from permafrost soils therefore depend on how and where the landscape undergoes abrupt thaw, and this spatial pattern remains one of the largest sources of uncertainty in the permafrost-carbon feedback \citep{walter201821st, turetsky2020carbon}.

Beyond climate feedbacks, abrupt thaw represents an immediate hazard to the roughly 5 million people who live in the Arctic circumpolar permafrost region \citep{ramage2021population}. Projections suggest that 30-50\% of Arctic infrastructure may be at risk of damage by permafrost thaw by 2050, threatening tens of billions of U.S. dollars in damage by the end of the century \citep{hjort2022impacts}. In Alaska alone, permafrost thaw may cause \$37 to \$51 billion in damage to buildings and roads \citep{manos2025permafrost}. This burden will fall disproportionately on remote and Indigenous communities, whose unofficial, ``vernacular'' infrastructure and assets often lie outside conventional damage inventories \citep{verstraete2025vernacular}. Between 1974 and 2019, thermokarst at Point Lay, Alaska, expanded nearly an order of magnitude faster beneath the developed village than in the surrounding tundra, indicating a potential positive feedback between the built environment and abrupt thaw \citep{jones2025climate}. Predicting how and where abrupt thaw will develop on the landscape is therefore critical for helping Arctic communities adapt to or mitigate the effects of future climate change.

Abrupt permafrost thaw proceeds faster than gradual thaw, by definition, but also results from different processes, leaves distinct surface expressions, and has specific consequences for ecosystems, hydrology, and carbon release \citep{webb2025review}. Gradual thaw typically proceeds from the top down as a progressive deepening of the seasonally thawed active layer into near-surface permafrost \citep{walter201821st, webb2025review}. Abrupt thaw, by contrast, is not confined to this vertically ordered progression: subsidence, erosion, mass wasting, and the flow of water and heat can rapidly propagate thaw downward and laterally, affecting deeper or adjacent permafrost (Fig. \ref{fig:process-schematic}) \citep{turetsky2020carbon, nitzbon2020fast, webb2025review}. The resulting changes to the landscape occur quickly relative to geomorphic time scales, from vertical subsidence of almost a meter over a decade \citep{farquharson2019climate} to thaw-slump headwall retreat of more than ten meters in a single year \citep{kokelj2015increased}. A site's thaw mode, abrupt or gradual, is not dictated by the presence of permafrost alone, but emerges from an interplay of ground ice, geomorphology, hydrology, and disturbance history \citep{jorgenson2005response, kokelj2013advances}.

Permafrost has been mapped extensively across the Arctic, but existing products largely resolve conditions other than the mode of thaw. Many characterize the substrate, outlining permafrost extent and zonation \citep{brown1997circum}, ground ice content \citep{jorgenson2008permafrost}, and probability of occurrence \citep{obu2019northern}. Others document where abrupt thaw features have formed, such as ponds and lakes \citep{muster2017perl} or retrogressive thaw slumps \citep{nitze2025darts}, or quantify how far degradation of specific features has advanced \citep{zhang2025mapping}. A growing set of susceptibility maps predicts where specific abrupt thaw processes may develop, including thermokarst lakes \citep{wang2023thermokarst}, retrogressive thaw slumps \citep{makopoulou2024retrogressive}, and active layer detachment \citep{makopoulou2025susceptibility}, and a recent hemispheric model combines lake and slump susceptibility into a single product \citep{yang2025future}. Each of these efforts maps the susceptibility of the landscape to a named landform, estimating the likelihood that a particular feature will form at a given point, and the combined model remains a union of predictions anchored to specific processes. What none of them map is thaw mode itself: whether a site's conditions favor the abrupt or gradual pathway, independent of any single landform (Fig. \ref{fig:process-schematic}). Across Alaska, that contrast is carried only by point-wise samples from expert analysis of field and remote sensing observations \citep{webb2026comprehensive}, and not yet mapped continuously.

Here, we extend this susceptibility-mapping approach to thaw mode, producing a continuous, predictive map of Alaska's permafrost region. We train a gradient-boosted classifier on 19,288 expert-labeled thaw points from the Alaska Permafrost Thaw Database (v2.0) \citep{webb2026comprehensive}. The classifier evaluates sites along a gradient of how strongly key environmental variables resemble either abrupt or non-abrupt thaw, using 70 geospatial indicators sourced from remote sensing, community data products, and model reanalysis. To mitigate sampling bias inherent to the Thaw Database, we evaluate the model under spatial cross-validation, withholding entire geographic regions from training but not testing, bounded by an applicability mask that indicates where points fall within the envelope of the training data, and then apply it to a statewide feature grid populated by the same data sources. The resulting product is a log-evidence index of thaw mode that is \textit{prior-free}: it characterizes the evidence from local features alone, without assuming how common either thaw mode is across the landscape. This index represents the relative susceptibility of the landscape to abrupt thaw, as compared to gradual thaw. We then use grouped Shapley (SHAP) values, which attribute the model output across its input features \citep{lundberg2017unified, lundberg2020local}, to identify which geospatial indicators best distinguish between thaw modes and how they shift the classifier's prediction. This paper describes the expert-labeled points and environmental predictors in Section 2, model construction and evaluation in Section 3, the thaw mode map and its drivers in Section 4, and limitations, interpretations, and comparisons with the broader permafrost mapping literature in Section 5.

\begin{figure}[hbtp]
	\centering
	\includegraphics[width=\textwidth]{figures/01_process_schematic.pdf}
	\caption{
    	\textbf{Conceptual illustration of permafrost thaw mode.}
    	Idealized cross-sections (a), not drawn to scale, contrasting the two thaw modes and the ground-ice characteristics of each. Ice-poor permafrost (top left) leads to non-abrupt thaw (top right), where the active layer may deepen but the surface is broadly maintained, allowing gradual subsidence at most. Ice-rich permafrost (bottom left) undergoes abrupt thaw (bottom right) when excess ground ice is removed, leading to pronounced surface reorganization. Field photographs show two examples of abrupt thaw landforms: a thaw slump along an eroding riverbank in Alaska's North Slope (b), photographed by Irina Overeem, and a drone image of a thermokarst wetland in Goldstream Valley (c), photographed by Kevin Rozmiarek.
	}
	\label{fig:process-schematic}
\end{figure}

\end{document}
